\section*{Chapitre \ref{ch:nw:fixed-income-and-rfq-platform-connectivity} --- Connectivité des plateformes de taux et RFQ}

\begin{solution}{exo:nw:fixed-income-and-rfq-platform-connectivity:1}
Il n'est pas montré au client~: le dealer n'a aucun prix dans cette enchère, quel qu'il aurait été.
\end{solution}

\begin{solution}{exo:nw:fixed-income-and-rfq-platform-connectivity:2}
Les marchés américains du CME dans le NY5 à Secaucus, avec reprise à Aurora~; les marchés européens dans le LD4.2 à Slough, avec reprise dans le NY5.
\end{solution}

\begin{solution}{exo:nw:fixed-income-and-rfq-platform-connectivity:3}
Les étapes automatiques totalisent quelques dizaines de millisecondes et n'atteignent pratiquement jamais deux secondes~; une personne, dont le temps médian est de six secondes, dépasse
deux secondes dans 98.6~\% des cas, donc la probabilité est d'environ $0.02 \times 0.986 = 1.97~\%$, soit 1.94~\% dans la simulation.
\end{solution}

\begin{solution}{exo:nw:fixed-income-and-rfq-platform-connectivity:4}
À l'échéance, le client compare toutes les réponses reçues à temps, quel que soit leur ordre~; les réponses des dealers arrivent déjà à temps sauf
celle de la personne, que ne modifie pas un calculateur de prix plus rapide.
\end{solution}

\begin{solution}{exo:nw:fixed-income-and-rfq-platform-connectivity:5}
Avec cinq dealers interrogés, la meilleure des cinq premières est la meilleure parmi tous ceux qui répondent à temps, sauf si quelqu'un répond après le temporisateur~: c'est le
«~meilleur à l'échéance~», à ceci près que l'exécution a lieu dès que la cinquième réponse arrive. Avec davantage de dealers interrogés, il devient «~meilleur parmi les
cinq premières~».
\end{solution}

\begin{solution}{exo:nw:fixed-income-and-rfq-platform-connectivity:6}
Les contrats à terme sur bons du Trésor se négocient à Aurora~; avec un carnet au comptant là aussi, les deux jambes d'un basis trade se trouvent dans un même bâtiment, et les millisecondes
du corridor New Jersey--Chicago ainsi que ses coûts disparaissent de la transaction.
\end{solution}

\begin{solution}{exo:nw:fixed-income-and-rfq-platform-connectivity:7}
À environ \qty{50}{\milli\second}, le temps de calcul propre aux concurrents~: 19.0~\% avec «~première acceptable~» contre 19.2~\% à l'échéance~; plus rapide
que les concurrents, «~première acceptable~» avantage le dealer~; plus lent, il le pénalise.
\end{solution}

\begin{solution}{exo:nw:fixed-income-and-rfq-platform-connectivity:8}
Le temporisateur n'est pas la course lorsque les clients exécutent la première réponse acceptable ou la meilleure parmi les premières~: le dealer court alors contre
ses concurrents, dont les calculateurs de prix prennent quelques dizaines de millisecondes, et dix millisecondes valent près de quinze centimes de prix.
\end{solution}

\begin{solution}{pb:nw:fixed-income-and-rfq-platform-connectivity:1}
\textbf{Partie I.}
\begin{enumerate}
\item 1.94~\%, presque entièrement issus des 2~\% de demandes traitées par une personne.
\item Zéro à deux secondes~; 2.4~\% avec un temporisateur de \qty{100}{\milli\second}.
\item Au-delà d'environ une demi-seconde~: la simulation trouve 0.002~\% à \qty{562}{\milli\second} et aucun au-delà d'une seconde.
\item Les deux ne donnent aucune transaction, mais une réponse manquée indique aussi au client que le dealer n'a pas répondu~: un point contre lui dans le
  classement des dealers du client.
\end{enumerate}
\textbf{Partie II.}
\begin{enumerate}[resume]
\item 19.2~\% à l'échéance, 26.4~\% pour la meilleure des trois premières, 39.0~\% pour la première acceptable.
\item 0, 1.6~\% et 5.7~\% des demandes.
\item À l'échéance, toutes les réponses reçues à temps sont comparées, de sorte que la meilleure réponse acceptable est toujours retenue.
\item Il est plus rapide que ses quatre rivaux (20 contre \qty{50}{\milli\second} de calcul), il fournit donc souvent la première réponse acceptable.
\end{enumerate}
\textbf{Partie III.}
\begin{enumerate}[resume]
\item 0.0, 2.3 et 12.6 points.
\item 0.9, 1.0 et 0.9 points.
\item Environ quinze (12.6 contre 0.9 points par centime).
\item À partir de son propre historique~: gains et pertes selon son temps de réponse et son rang d'arrivée, client par client~; un taux de gain qui baisse avec le
  rang, à prix égal, révèle une règle autre que l'échéance.
\end{enumerate}
\textbf{Partie IV.}
\begin{enumerate}[resume]
\item \emph{Résultat nommé}~: avec un pipeline dont les étapes automatiques prennent quelques dizaines de millisecondes et 2~\% des demandes traitées par une personne,
  le dealer manque 1.94~\% des temporisateurs de deux secondes~; sous une règle de première acceptable, il perd 5.7~\% de toutes les demandes au profit de dealers plus rapides, malgré
  le meilleur prix acceptable, et diviser par deux son temps de calcul rapporte 12.6 points de taux de gain contre 0.9 pour un centime de prix.
\item Publié~: les protocoles et l'automatisation des plateformes, les sites des carnets d'ordres. Supposé~: tous les chiffres du pipeline, la part de la personne et
  ses délais, le temporisateur, les paramètres du modèle de prix et l'usage des règles.
\item Sur l'étape de calcul si ses clients utilisent des règles de première acceptable ou de premier-$k$~; sur la part envoyée à une personne si les temporisateurs sont courts~; sur
  le prix s'ils attendent l'échéance.
\item L'avantage disparaît et les taux de gain reviennent vers un sur cinq~; la rapidité protège alors la part de marché plutôt que de la conquérir.
\item Le réduire en élargissant les limites de l'auto-répondeur là où le risque le permet, et acheminer le reste vers une personne avec une alerte et une
  échéance bien en deçà du temporisateur.
\item Tout~: un carnet d'ordres est une bourse, avec flux, passerelles et \omterm{def:nw:colocation-products-and-how-they-are-sold:colo}{colocation}, et les microsecondes redeviennent importantes.
\item Dans les lignes de taux~: les plateformes et les sessions, le budget de latence de l'auto-répondeur, la file d'attente de la personne, et les sites des
  carnets d'ordres et leurs \omterm{def:nw:colocation-products-and-how-they-are-sold:mmr}{interconnexions} (cross-connects).
\item Lorsque la règle du client s'exécute avant le temporisateur, sur la première réponse acceptable ou les toutes premières réponses.
\end{enumerate}
\end{solution}

\begin{solution}{iq:nw:fixed-income-and-rfq-platform-connectivity:1}
Réseau en entrée, moteur FIX, enrichissement de la demande, calcul du prix (fair value, courbe, comparables, inventaire), limites et risque, réponse en sortie. Le calcul,
s'il recalcule une courbe ou attend des données~; une personne, si la demande sort des limites automatiques.
\iqlookfor{Étapes~; la variance du calculateur~; le chemin manuel.}
\end{solution}

\begin{solution}{iq:nw:fixed-income-and-rfq-platform-connectivity:2}
Pour exécuter vite pendant que le marché est encore là où il l'a vu, pour réduire la fuite d'information d'une longue enchère, et parce que ses propres
règles ou sa charge de travail imposent l'exécution automatique dès que le prix est assez bon.
\iqlookfor{Vitesse d'exécution~; fuite d'information~; automatisation~; tolérance de prix.}
\end{solution}

\begin{solution}{iq:nw:fixed-income-and-rfq-platform-connectivity:3}
On journalise l'arrivée de chaque demande, le temps de réponse, le rang d'arrivée si la plateforme le communique, le prix par rapport au prix gagnant ou au prix
de couverture (cover price), et l'issue~; on modélise la probabilité de gain en fonction du prix et du temps de réponse conjointement, client par client.
\iqlookfor{Modèle conjoint du prix et du temps~; règles par client~; le prix de couverture.}
\end{solution}

\begin{solution}{iq:nw:fixed-income-and-rfq-platform-connectivity:4}
La latence par étape pour localiser l'étape~; les dépendances de données du calculateur, une reconstruction lente de la courbe ou un défaut de cache (cache miss), le ramasse-miettes ou
l'allocation mémoire, la journalisation, les threads du moteur FIX, et si la part envoyée à une personne a changé.
\iqlookfor{Attribution par étape~; dépendances de données~; effets du runtime~; la part manuelle.}
\end{solution}

\begin{solution}{iq:nw:fixed-income-and-rfq-platform-connectivity:5}
Cela dépend des règles d'exécution des clients~: on estime le taux de gain marginal d'un centime de prix et d'une milliseconde à partir de l'historique propre du dealer, puis on
les pèse contre leurs coûts~: un centime est cédé à chaque transaction, une milliseconde ne se paie qu'une fois.
\iqlookfor{Les règles décident~; valeurs marginales~; coûts récurrents contre coûts ponctuels.}
\end{solution}

\begin{solution}{iq:nw:fixed-income-and-rfq-platform-connectivity:6}
Sessions FIX vers chaque plateforme avec des lignes redondantes~; un service de calcul de prix unique avec un budget de latence et un chemin rapide pour les obligations liquides~; limites
et risque en ligne~; une file d'attente de personnes avec alertes~; une surveillance par plateforme des échecs, des rangs et des pertes face à la vitesse~; \omterm{def:nw:colocation-products-and-how-they-are-sold:colo}{colocation} et
\omterm{def:nw:colocation-products-and-how-they-are-sold:mmr}{interconnexions} (cross-connects) pour le carnet d'ordres dans son centre de données, avec ses propres feed handlers et sa passerelle.
\iqlookfor{Sessions~; budget du pipeline~; chemin manuel~; surveillance~; carnet d'ordres comme bourse.}
\end{solution}
